\section{データベース設計}
本システムのデータベース設計を以下に示す．図 \ref{fig:sec8/er} はデータの関係をER図で表したものであり，表 \ref{tb:user} から表　はデータベースのテーブル設計である．
\begin{figure}[H]
    \centering
    \includesvg[scale=0.8]{sections/section8_database.svg}
    \caption{ER図}
    \label{fig:sec8/er}
\end{figure}

\if0
\begin{table}[H]
    \centering
    \caption{}
    \label{tb:}
    \begin{tabular}{|c|c|c|c|c|c|} \hline
        カラム名 & PK & FK & その他の制約 & 論理データ型 & 説明 \\ \hline\hline      
    \end{tabular}
\end{table}
\fi

\begin{table}[H]
    \centering
    \caption{会員情報}
    \label{tb:user}
    \begin{tabular}{|c|c|c|c|c|c|} \hline
        カラム名 & PK & FK & その他の制約 & 論理データ型 & 説明 \\ \hline\hline
        id & ○ & & not null & VARCHAR(50) & Google ID \\ \hline
        gmail & & & not null & VARCHAR(100) & メールアドレス \\ \hline
        role & & & not null & ENUM('user','business','admin') & 会員区分 \\ \hline
    \end{tabular}
\end{table}